\begin{idstatement}{D.1.15}{Proved comparison with identical integral monodromy}
\label{st:D.1.15}
Two smooth proper families of two-tori over $\PP^1$ have the same
constant marked $H^1(-;\Z)$ local system but have
$R^1f_*\OO=\OO^2$ and $R^1f_*\OO=\OO(1)^2$, respectively.
Thus the integral representation alone does not determine the
coherent splitting for general compact torus fibrations. Both
examples have $b_1=4$; this comparison does not assert a new
counterexample under $b_1=b_2=0$.
\end{idstatement}
\begin{proof}
Let $V=\OO_{\PP^1}(1)^{\oplus2}$. On the affine coordinate $t$,
four global sections of $V$ are the columns of
\[
P(t)=\begin{pmatrix}1&i&t&it\\-t&it&1&-i\end{pmatrix}.
\]
They are a real basis in every fibre. Indeed, writing
$z_1=a+ib$, $z_2=c+id$, their real linear combination is
$(z_1+t z_2,\overline z_2-t\overline z_1)$.
Its vanishing gives $(1+|t|^2)\overline z_2=0$ and then $z_1=0$.
At infinity, the $\OO(1)$ frame changes $P(t)$ to $t^{-1}P(t)$,
whose leading columns again form a real basis. The lattice
generated by these sections therefore acts freely and properly by
holomorphic fibre translations on $V$. Its quotient $Z$ is a
compact smooth holomorphic torus bundle over $\PP^1$.
The displayed real basis identifies it differentiably with
$\PP^1\times(\R^4/\Z^4)$ and marks its constant integral local
system.

The fibrewise period map includes the bundle of holomorphic
one-forms by $P(t)^t$. The Dolbeault coefficient map gives the
exact sequence
\[
0\longrightarrow\OO(-1)^2
 \xrightarrow{P(t)^t}\OO^4
 \longrightarrow R^1f_*\OO_Z\longrightarrow0.
\]
The last map can be represented explicitly by
\[
B(t)=\begin{pmatrix}-it&-t&i&1\\-i&1&-it&t\end{pmatrix}.
\]
The identity $B(t)P(t)^t=0$ holds polynomially. Its two rows have
rank two at every finite $t$, as is seen from their two-by-two
minors, and their leading coefficients have rank two at infinity.
Thus this is the exact complex of vector bundles
\[
0\longrightarrow\OO(-1)^2\longrightarrow\OO^4
 \xrightarrow{B}\OO(1)^2\longrightarrow0,
\]
proving $R^1f_*\OO_Z\simeq\OO(1)^2$. The polynomial identity is
also checked in the attached exact program.

For the constant family $Z_0=Z_{t=0}\times\PP^1$, the same
marked fibre lattice gives the identical constant integral local
system, whereas coherent K\"unneth gives
$R^1f_*\OO_{Z_0}=H^1(Z_{t=0},\OO)\otimes\OO=\OO^2$.
The different splittings are therefore compared by an explicit
identity of the integral representations and explicit coherent
quotient maps.
\end{proof}


\subsection{D.2. The trivial-monodromy question}

% D.2.1 from part_D.tex

